\documentclass[11pt]{ctexart}
\usepackage{amsmath,amssymb,longtable}
\begin{document}
\section*{P151: 当前完整数学模型（2026-10-20）}
\section*{修订依据与边界}恢复源文第6、11、27–28页客流/库存含义，使用完整累计到达积分及前序候车递推；当前source_horizon_branch仅原有限时域，无额外名义发车窗。全部10例已有新分支真实最优及独立验证。

\section*{集合与决策}
每列车$k$有完整有序区段集合$E_k$及有序车站集$U_k$。
编组大小$n\in\mathcal N$，区段耗时$D_{ke}(n)$，共享资源、OD到达率$r^{uv}_p$、
周期长度$\Delta$、车辆库存$F_d$、耦合$T_c$、解编$T_d$、折返$\hat T$、
停站$T_w$及有限时域$H$均由实例给出。
$z_{kn}\in\{0,1\}$、$s_k=\sum_nn z_{kn}$、$t_{ke}\in\{1,\ldots,H\}$。
\[
\sum_nz_{kn}=1,\quad D_{ke}=\sum_nD_{ke}(n)z_{kn},\quad
t_{k,e+1}=t_{ke}+D_{ke},\quad t_{ke}+D_{ke}\le H.
\]
每对占用同一资源的不同列车区段，二元$o$选择唯一先后：
\[
t_{jf}\ge t_{ie}+D_{ie}+h-M(1-o),\quad
t_{ie}\ge t_{jf}+D_{jf}+h-Mo.
\]
\section*{库存与周转}
同库发车列车$i,j$的二元$\omega_{ij}$满足
$t_{j0}\ge t_{i0}+1-M(1-\omega_{ij})$、
$t_{i0}\ge t_{j0}+1-M\omega_{ij}$，$\omega_{ji}=1-\omega_{ij}$。
对到达该库$j$和离开该库$i$令
$g_{ij}=t_{j,last}+T_d(s_j-1)+\hat T+T_c(s_i-1)-t_{i0}$。
\[
-My_{ij}\le g_{ij}\le M(1-y_{ij}),\qquad
s_i\le F_d-\sum_{j\ne i}\omega_{ji}s_j+\sum_{j\text{到库}}y_{ij}s_j.
\]
$y$及$\omega$与编组二元的乘积采用$v\le y,v\le z,v\ge y+z-1,v\ge0$精确线性化。
等时返回可用库存；若$g=0$可选$y=0$是保守可行选择，不产生额外车辆。
\section*{时变客流及完整目标}
列车的车站时间是对应区段起点时间。定义完整累计到达量
\[
A_{uv}(t)=\sum_p r^{uv}_p\min\{\Delta,\max(0,t-p\Delta)\}.
\]
通过二元周期选择$a_{kup}$及整数周期内偏移$c_{kup}$精确表示：
$\sum_pa=1$，$0\le c_{kup}\le\Delta a_{kup}$，
$t_{ku}+T_w=\sum_p(p\Delta a_{kup}+c_{kup})$，
$A_{uv}(t_{ku}+T_w)=\sum_p[(\sum_{q<p}r_q^{uv})\Delta a_{kup}+r_p^{uv}c_{kup}]$。
每个OD的列车顺序由输入服务标识给定，并约束相邻服务$t_{ku}\ge t_{ju}+1$。
\[
W_{1uv}=A_{uv}(t_{1u}+T_w),\quad
W_{kuv}=W_{juv}-B_{juv}+A_{uv}(t_{ku}+T_w)-A_{uv}(t_{ju}+T_w),
\quad0\le B_{kuv}\le W_{kuv}.
\]
\[
C_{ku}=c_0s_k-\sum_{a\prec u}\sum_{v\succ u}B_{kav}\ge0,\quad
\sum_{v\succ u}B_{kuv}\le C_{ku},\quad
P^{out}_{ku}=\sum_{a\prec u}B_{kau}.
\]
\[
\min\sum_{k,u,v}(W_{kuv}-B_{kuv}).
\]
名义时刻仅提供MIP初值和服务编号；不添加名义发车窗口。完整当前模型对应source\_horizon\_branch。

\end{document}
